\documentclass{beamer}

\usepackage[BoldFont,SlantFont]{xeCJK}
\setCJKmainfont[Path=fonts/]{STKaiti.ttf}%如使用overleaf编译此处报错，则点击页面菜单将编译器更改为xelatex即可
\setCJKsansfont[Path=fonts/]{STKaiti.ttf}
\setCJKmonofont[Path=fonts/]{STKaiti.ttf}
\setmainfont{Times New Roman}
\setsansfont{Times New Roman}
\setmonofont{Times New Roman}

\usepackage{fontenc}
\usepackage{hyperref}
\usepackage{ragged2e}
\usepackage{lipsum}
\usepackage{setspace}
\usepackage[justification=centering]{caption}

\usepackage{graphicx}
\usepackage{subfig}

% other packages
\usepackage{latexsym,amsmath,xcolor,multicol,booktabs,calligra}
\usepackage{graphicx,pstricks,listings,stackengine}

\usepackage{bm}
\usepackage{amssymb}
\usepackage{multirow}
\usefonttheme{serif}
\setbeamertemplate{navigation symbols}{}

\author{第一作者 \inst{1}，第二作者 \inst{2*}}
\title{论文标题}
\subtitle{论文副标题}
\institute[*]{\inst{1} 广东工业大学 \ \ 部门(院), 1431***78@qq.com \\
\inst{2*} 二作单位 \ \ 二作部门, wh****8@126.com}
\date{2023年11月25日}
\usepackage{gdut}

% defs
\def\cmd#1{\texttt{\color{red}\footnotesize $\backslash$#1}}
\def\env#1{\texttt{\color{blue}\footnotesize #1}}
\definecolor{deepblue}{rgb}{0,0,0.5}
\definecolor{deepred}{rgb}{0.6,0,0}
\definecolor{deepgreen}{rgb}{0,0.5,0}
\definecolor{halfgray}{gray}{0.55}

\lstset{
    basicstyle=\ttfamily\small,
    keywordstyle=\bfseries\color{deepblue},
    emphstyle=\ttfamily\color{deepred},    % Custom highlighting style
    stringstyle=\color{deepgreen},
    numbers=left,
    numberstyle=\small\color{halfgray},
    rulesepcolor=\color{red!20!green!20!blue!20},
    frame=shadowbox,
}

\setbeamertemplate{caption}[numbered]
\renewcommand\figurename{图} 
\renewcommand\tablename{表}

\usepackage{perpage} %the perpage package
\MakePerPage{footnote} %the perpage package command

\newfontfamily\wingdingfamily[Path=fonts/]{wingding.ttf}
\newcommand\headtoright{%
  \texorpdfstring{{\wingdingfamily\XeTeXglyph190\relax}}{???}
}

\setbeamertemplate{itemize item}{\headtoright}
\setbeamertemplate{itemize subitem}{$\bullet$}
\setbeamertemplate{itemize subsubitem}{$\circ$}



%%%%%%%%%%%%%%%%%%%%%%%%本PPT模板修改自苏州大学模板，由广工管院21级学生卢义桢上传。科研不易，希望为诸位学术之路略尽绵薄之力！%%%%%%%%%%%%%%%%%%%%%%%%%%

\begin{document}
\begin{frame}
    \titlepage
    \begin{figure}[htpb]
        \begin{center}
            \includegraphics[width=0.19\linewidth]{figs/GDUT_logo.eps}
        \end{center}
    \end{figure}
\end{frame}

\begin{frame}
    \tableofcontents[sectionstyle=show,subsectionstyle=show/shaded/hide,subsubsectionstyle=show/shaded/hide]
\end{frame}


\section{选题背景}             %PPT的章节，可以自由设置
\begin{frame}{电*****问题}     %第一章节的第一页PPT，每一个frame环境为一页PPT
	\begin{columns}[c]
		\column{6.5cm}
		\begin{itemize}
			\item UPS 2019年的调查报告指出，*****\textcolor{red}{1/3}的网购者会巴拉巴拉；      %强调数据部分标红
			\item 巴拉巴拉；
			\item 巴拉巴拉。
		\end{itemize}
		\bigskip
		\bigskip	
		\column{5cm}
		\begin{figure}[htbp]
			\centering       
			\includegraphics[width=0.75\textwidth]{figs/return.eps} 
			                  %插入图片也可用png/jpg格式，但推荐.eps格式
			\caption{图名 \\ 
				\tiny（数据来源：UPS Pulse of The Online Shopper 2019）}  
		    % 中英文的逗号、分号、括号等都可以，但为了美观建议全文统一为中文或英文符号。 
			\label{123}        %引用本图的标签
		\end{figure}
	\end{columns}
\end{frame}
\begin{frame}{巴拉巴拉的应用}   %第一章节第二张PPT
	 \begin{columns}[c]        %页面左右分栏
		\column{5cm}
		增强现实技术:\\
		\vspace{0.1in}         %空行
		\footnotesize  巴拉巴拉\\
		\vspace{0.2in}
		\begin{itemize}
			\item Deloitte Digital公司调查数据显示：巴拉巴拉服务可将线上消费的\textcolor{red}{降低25\%} \footnote{\href{https://www.gdut.edu.cn/}{数据来源网址链接}}; 
			                   %增加脚注：呐，这就叫专业！
			\item 巴拉巴拉。
		\end{itemize}
		\bigskip               %增加空行，控制行间距。
		\bigskip
		\column{5cm}           %自由设置column宽度，直到左右栏大小调整到位
		\begin{figure}[htbp]
			\centering       
			\subfloat
			{
				\label{fig:subfig1}\includegraphics[width=0.2\textwidth]{figs/GDUT666.eps}
			}
			\subfloat
			{
				\label{fig:subfig2}\includegraphics[width=0.2\textwidth]{figs/GDUT666.eps}
			}
			\caption{图名}
			\label{fig:subfig_1}
		\end{figure}
	\end{columns}
\end{frame}
\begin{frame}{巴拉巴拉的问题}
	\begin{minipage}{0.9\linewidth}   %页面上下分栏
		\begin{table}[htbp]           %三线表
			\centering
			\caption{不同平台巴拉巴拉情况}
			\scalebox{0.6}{
				\begin{tabular}{c c l}
					\toprule
					企业 & 销售模式 & 应用现状 \\
					\midrule
					巴啦 & 代理模式 & 巴拉巴拉国内领先巴拉巴拉巴拉巴拉巴拉，\\
					& & 自称全球规模最大精度最高的服务。\\
					巴拉 & 转售模式+代理模式 & 平台巴拉巴拉服务。\\
					东 & 转售模式+代理模式 & 巴拉巴拉等。\\
					拼 & 代理模式 & 巴拉巴\\
					...&...&... \\
					\bottomrule
			\end{tabular}}
		\end{table}
	\end{minipage}\hspace{1cm}
	\bigskip
	\begin{minipage}{0.9\linewidth}
			\begin{figure}[htbp]     
			\centering       
			\subfloat                   %并排多张图片
			{
				\label{fig:subfig3}\includegraphics[width=0.287\textwidth]{figs/bg4.eps}
			}
			\subfloat
			{
				\label{fig:subfig4}\includegraphics[width=0.251\textwidth]{figs/bg3.eps}
			}
			\subfloat
			{
				\label{fig:subfig5}\includegraphics[width=0.261\textwidth]{figs/bg2.eps}
			}
			\subfloat
			{
				\label{fig:subfig6}\includegraphics[width=0.25\textwidth]{figs/bg1.eps}
			}
			\caption{巴拉巴拉数量}
			\label{fig:subfig_2}
		\end{figure}
	\end{minipage}
\end{frame}

\section{模型设置}
\subsection{问题描述}
\begin{frame}{问题描述}
	\justifying{
    考虑一个巴拉巴拉（$d$）、电巴拉巴拉（$v$）与巴拉巴拉链，巴拉巴拉与巴拉巴拉之间存在两种巴拉巴拉，分别为转巴拉巴拉（$y$）和代巴拉巴拉（$i$） 巴拉巴拉的影响。}
    \begin{figure}[t!]
    	\flushleft       %居左布置
    	\subfloat{        
    	\includegraphics[width=0.1\textwidth]{figs/GDUT666.eps}}
    	\flushright      %居右布置
    	\subfloat{
    	\includegraphics[scale=0.011]{figs/GDUT6666.eps}}
    %图片的布置需要反复调整大小到合适为止，scale与width=**\textwidth均是用于控制图片大小的命令。一般只用Latex作表，作图则用Visio、Origin或Matlab画好插入就完事了
    \caption{巴拉巴拉巴拉巴拉序}
    \end{figure}
\end{frame}
\subsection{参数设置及假设}
\begin{frame}{参数设置及假设}
	\begin{columns}
		\column{6cm}
		\begin{table}[htbp]
			\centering
			\caption{模型参数符合及含义}
			\scalebox{0.65}{
				\begin{tabular}{c c}
					\toprule
					符号 & 含义 \\
					\midrule
					$i$ & 不同模式, $i\in \{R,M\}$, $R$:巴拉式, $M$：巴拉巴拉巴拉 \\
					$j$ & 不同主体, $j\in\{m,e\}$, $m$:制巴拉商, $e$: 巴拉商巴拉 \\
					$\Pi^{i}_j$ & $i$情境下$j$的巴拉润 \\
					$D$ & 产品*** \\
					$p$ & 销售*** \\
					$w$ & ***格 \\
					$c$ & 单位巴拉巴拉 \\
					$q$ & 产品*** \\
					$\lambda$ & ***率 \\
					$r$ & 退款 \\
					$a$ & 产品 \\
					$b$ & 巴拉巴拉的敏感性系数 \\
					$s$ & ***值 \\
					$F$ & ***成本 \\
					$\theta$ & 巴拉巴拉的影响 \\
					$\xi_0$ & 巴拉巴拉率 \\
					\bottomrule
				\end{tabular}}
		\end{table}
	\bigskip
	\bigskip
	\column{5cm}
	\begin{itemize}
		\tiny
		\item 参考以往学者的做法，定义需求为：$D=a-bq$，巴拉巴拉为1；
		\item 产品的巴拉巴拉的应用有关，令$\xi_{(q)}=(1-\theta)\xi_{0}-q$；
		\item 提高产品质量带来的成本巴拉巴拉关系，即$c_{(q)}=q^{2}/2$；
		\item 应用巴拉巴拉；
		\item 为保证商家巴拉巴拉
		\item 由于巴$r>s$。
	\end{itemize}
	\end{columns}
\end{frame}

\section{模型求解}
\subsection{巴拉巴拉}
\begin{frame}{巴拉巴拉式}
	\begin{columns}
		\column{6cm}
		\begin{exampleblock}{巴拉巴拉利润}
			\begin{small}
				\begin{equation*}
					\begin{split}
					\bm{\Pi^{R}_{m(w,q)}}&=(2)D-c_{3}\\
					&=(p)(4)-\frac{ttt}{ttty}
					\end{split}
				\end{equation*}
				\begin{equation*}
					\begin{split}
						\bm{\Pi^{R}_{e(p)}}&=(qw)D-\xi_{er}(4)D-yy\\
						&=[p456\xi_{(q)}-\theta\xi_{x}](9)
					\end{split}
				\end{equation*}
			\end{small}
		\end{exampleblock}
		\column{5cm}
		\begin{figure}
			\centering       
			\includegraphics[width=0.32\textwidth]{figs/GDUT666.eps}
			\caption{模型结构}
		\end{figure}
	\end{columns}
\end{frame}
\subsection{巴拉巴拉模式}
\begin{frame}{代巴拉巴拉模式}
	\begin{columns}
		\column{7cm}
		\begin{exampleblock}{制巴拉巴拉利润}
			\begin{footnotesize}
				\begin{equation*}
					\begin{split}
						\bm{\Pi^{M}_{m(p,q)}}&=(fg)(t)D-(r-s)\xi_{(4)}67-c_{(aasd)}-ertASF\\
						&=[(456-\lambda qwe)(wer-c)-(r-s)\xi_{7}-\theta \xi_{0}](v)-\frac{v^{2}}{m}
					\end{split}
				\end{equation*}
				\begin{equation*}
					\begin{split}
						\bm{\Pi^{M}_{e}}=\lambda 45rtt=fgt(tg)89j
					\end{split}
				\end{equation*}
			\end{footnotesize}
		\end{exampleblock}
		\column{4cm}
		\begin{figure}
			\centering       
			\includegraphics[width=0.32\textwidth]{figs/GDUT666.eps}
			\caption{模型结构}
		\end{figure}
	\end{columns}
\end{frame}
\begin{frame}{均衡巴拉巴拉}
    \begin{exampleblock}{决策变量巴拉巴拉 \scriptsize{（此时balabal需要满足条件：$0<b<\frac{erd}{(r-s)^2}$）}}
    	\begin{scriptsize}
    		\begin{equation*}
    			\centering
    			\begin{split}
					\bm{\Pi^{M}_{m(p,q)}}&=(fgddd)(t)D-(r-s)\xi_{(4)}67-c_{(aasd)}-ertdsASF\\
					&=[(456-\lambda qwe)(ssswer-c)-(r22e-s)\xi_{7}-\theta \xi_{0}](v)-\frac{v^{2}}{m}dsds
				\end{split}
			\end{equation*}
		\end{scriptsize}
	\end{exampleblock}
    \end{frame}

\section{结论启示}
\subsection{转巴拉巴拉下的相关结论}
\begin{frame}{巴拉巴拉格的影响}
	\begin{large}
		\flushleft
		推论一：\\
		\vspace{0.2in}
	\end{large}
    \justifying
    (1) $\frac{\partial D^{R*}}{\partial \theta}<0$，即巴拉巴拉低；\\
    \vspace{0.1in}
	(2) $\frac{\partial w^{R*}}{\partial \theta}>0$，即巴拉巴拉高；\\
	\vspace{0.1in}
    (3) $\frac{\partial q^{R*}}{\partial \theta}<0$，即\textcolor{red}{技巴拉巴拉低}；\\
    \vspace{0.1in}
    (4) $\frac{\partial p^{R*}}{\partial \theta}>0$，即巴拉巴拉越高。
\end{frame}
\begin{frame}{转售巴拉巴拉的影响}
	\begin{large}
		\flushleft
		对巴拉巴拉响：\\
		\vspace{0.1in}
	\end{large}
	\justifying 当$0<\xi_{0}<min(1,\frac{1-c}{ererd\theta)})$时，则$\frac{\partial \Pi^{R}_{m*}}{\partial frt}<0$；\\当$max(0,\frac{1-c}{4rtf})<\xi_{0}<1$时，则$\frac{\partial \Pi^{R*}_{m}}{\partial cv}>0$。即当巴拉巴拉较低时，制造商利润与巴拉巴拉呈负相关，反之巴拉巴拉。\\
	\vspace{0.2in}
	\begin{large}
		\flushleft
		对巴拉巴拉的影响：\\
		\vspace{0.1in}
	\end{large}
	\justifying 当$0<\xi_{0}<min(1,rf)$时，则$\frac{\partial \Pi^{R*}_{e}}{\partial der}<0$；\\当$max(0,\frac{1-c}{4r5t})<\xi_{0}<1$时，则$\frac{\partial \Pi^{R*}_{e}}{\partial dere}>0$。即巴拉巴拉。
\end{frame}

\subsection{巴拉巴拉的相关结论}
\begin{frame}{代巴拉巴拉的影响}
	\begin{large}
		\flushleft
		推论二：\\
		\vspace{0.2in}
	\end{large}
	\justifying
	(1) $\frac{\partial Dw^{err*}}{\partial 09898}<0$，即巴拉巴拉越低；\\
	\vspace{0.1in}
	(2) $\frac{\partial qww^{r4*}}{\partial advc}<0$，即\textcolor{red}{巴拉巴拉低}；\\
	\vspace{0.1in}
	(3) $\frac{\partial per^{ww*}}{\partial rrtr}>0$，即巴拉巴拉高。
\end{frame}

\subsection{不同模式的巴拉巴拉对比}
\begin{frame}{不同模式下巴拉巴拉比较}
	\begin{large}
		\flushleft
		命题一：\\
		\vspace{0.2in}
	\end{large}
	\justifying
	(1) 巴拉巴拉比较：当$0<\xi_{0}<min(1,tudhudh)$时，则$dsfsdf$，反之$dsder$。 \\
	\vspace{0.1in}
	(2) 巴拉巴拉比较：当$0<\xi_{0}<min(1,trtrtrt)$时，则$q^{r*}>q^{n*}$，反之$q^{tytyt*}<q^{uuyy*}$。\\
	\vspace{0.1in}
	(3) 巴拉巴拉比较：当$0<\xi_{0}<min(1,\frac{67dhjd}{2})$时，则$p^{mn*}>p^{ty*}$，反之$p^{rtt*}<p^{5ty*}$。
\end{frame}
\begin{frame}{不同模式下巴拉巴拉对比}
	\begin{large}
		\flushleft
		命题二：\\
		\vspace{0.2in}
		\justifying
		\begin{spacing}{1.5}
		$\frac{\partial q^{M*}}{\partial \theta}$，即巴拉巴拉巴拉巴拉巴拉巴拉巴拉巴拉巴拉巴拉巴拉巴拉巴拉巴拉。
		\end{spacing}
	\end{large}
\end{frame}
\subsection{巴拉巴拉的影响}
\begin{frame}{刻画消费巴拉巴拉}
	\begin{small}
		\justifying
		\begin{spacing}{1.5}
			根据以往文献的做法，消费者效用为需求的二次函数，因此$U_{(D)}=\frac{D}{b}-\frac{D^{2}}{2b}$，巴拉巴拉表达式如下：
		\end{spacing}
		\begin{equation*}
			CS=U(D)-pD=\frac{rffff}{b}-\frac{5rtt}{256555}
		\end{equation*}
	\bigskip
		\begin{equation*}
			SW=CS+\Pi_{e}+\Pi_{m}
		\end{equation*}
	\end{small}
\end{frame}

\subsection{数值分析}
\begin{frame}{数值分析}
	\begin{tiny}
		\begin{table}[htbp]  % 略微复杂的表格
			\centering
			\caption{巴拉巴拉对不同模式下均衡结果的影响
			\footnotesize{($\lambda=321,b=122,\xi=4,c=36,s=72,r=87$)}}
			\scalebox{0.9}{
				\begin{tabular}{c c c c c c c c c c}
					\toprule
					$\theta$&销售模式&$D^{*}$&$q^{*}$&$p^{*}$&$w^{R*}$&$\Pi^{*}_{e}$&$\Pi^{*}_{m}$&$CS^{*}$&$SW^{*}$ \\
					\midrule
					\multirow{2}{*}{0.1}&销售模式($R$)&0.2799&0.0460&3.8570&2.5986&0.2642&0.5273&1.0228&1.8251 \\
					&代理模式($M$)&0.4566&0.0913&2.7170&-&0.3311&0.7256&1.7705&2.8408\\
					\multirow{2}{*}{0.2}&销售模式($R$)&0.2283&0.0457&3.8587&2.5826&0.2605&0.5200&1.0159&1.8788 \\
					&代理模式($M$)&0.4520&0.0904&2.7399&-&0.3678&0.7110&1.7565&2.8659\\
					\multirow{2}{*}{0.3}&销售模式($R$)&0.2267&0.0453&3.8667&2.2165&0.2569&0.5127&1.6789&1.7876 \\
					&代理模式($M$)&0.4474&0.0895&2.7629&-&0.3306&0.6966&1.7424&2.7793\\
					\multirow{2}{*}{0.4}&销售模式($R$)&0.2251&0.0450&3.8747&2.5505&0.2532&0.5055&1.0020&1.7673 \\
					&代理模式($M$)&0.4428&0.0886&2.7859&-&0.3302&0.6824&1.7281&2.7486\\
					\multirow{2}{*}{0.5}&销售模式($R$)&0.2234&0.0447&3.8828&2.5345&0.2496&0.4983&0.9950&1.7481 \\
					&代理模式($M$)&0.4322&0.0876&2.8089&-&0.3298&0.6683&1.7138&2.7179\\
					\multirow{2}{*}{0.6}&销售模式($R$)&0.2218&0.0444&3.8908&2.5184&0.2461&0.4912&0.9880&1.7291 \\
					&代理模式($M$)&0.4336&0.0867&2.8319&-&0.3294&0.6543&1.6992&2.6571\\
					\multirow{2}{*}{0.7}&销售模式($R$)&0.4502&0.0440&3.8988&2.5024&0.2115&0.4841&0.9810&1.2201 \\
					&代理模式($M$)&0.4290&0.0858&2.8549&-&0.3288&0.6405&1.6992&2.6571\\
					\multirow{2}{*}{0.8}&销售模式($R$)&0.2186&0.0437&3.9068&2.4864&0.2390&0.4771&0.2340&1.6912 \\
					&代理模式($M$)&0.4244&0.0849&2.8779&-&0.3282&0.6269&1.6698&2.6255\\
					\multirow{2}{*}{0.9}&销售模式($R$)&0.2170&0.0434&3.9148&2.4703&0.2355&0.4701&0.9669&1.6724 \\
					&代理模式($M$)&0.4198&0.0453&2.9009&-&0.3276&0.6134&1.6549&2.5947\\
					\multirow{2}{*}{1.0}&销售模式($R$)&0.2154&0.0431&3.9228&2.4543&0.2321&0.4632&0.9599&1.6537 \\
					&代理模式($M$)&0.4152&0.0810&2.9239&-&0.1169&0.6090&1.6399&2.5009\\
					\bottomrule
			\end{tabular}}
		\end{table}
	\end{tiny}
\end{frame}
\begin{frame}{巴拉巴拉影响}
	\begin{figure}[htbp]
		\centering       
		\subfloat[$\lambda=0.2$]
		{
			\includegraphics[width=0.35\textwidth]{figs/fig1.eps}
		}
		\subfloat[$\lambda=0.4$]
		{
			\includegraphics[width=0.35\textwidth]{figs/fig2.eps}
		}
		\subfloat[$\lambda=0.6$]
		{
			\includegraphics[width=0.35\textwidth]{figs/fig3.eps}
		}
		\caption{巴拉巴拉的影响 ($\Delta \Pi^{*}_{e}=\Pi^{M*}_{e}-\Pi^{R*}_{e},\Delta \Pi^{*}_{m}=\Pi^{M*}_{m}-\Pi^{R*}_{m}$)}
	\end{figure}
\end{frame}

\section{参考文献}
\begin{frame}{参考文献}
	\tiny
%	\addcontentsline{toc}{section}{参考文献}
%	\nocite{1}             %只加入到参考文献列表中，不在文中引用（显示[id]）
    \nocite{*}             %显示所有文献
	\bibliography{ref}     %引用文献将.bib文件拷入ref.bib文件中
	\bibliographystyle{Style2014}   %控制文献显示风格
\end{frame}

\begin{frame}
    \begin{center}
        {\Huge\calligra Thanks!}
    \end{center}
\end{frame}

\end{document}